\documentclass[12pt]{article}
\usepackage[utf8]{inputenc}
\usepackage[T1]{fontenc}
\usepackage{hyperref}
\title{Kereho Language Resources for K-12 Curriculum Development}
\author{Compiled by Claude Sonnet 5.5 (claude-sonnet-5-5) for Andr\'as Kornai}
\date{}
\begin{document}
\maketitle

\section{Introduction}
Kereho (ISO 639-3 xke; Glottolog kere1285) is an Austronesian language of Borneo with about 500 speakers. Sources conflict on its identity. Wikipedia calls it Krio Dayak, a Kayan-Murik language of West Kalimantan. Joshua Project and Global Recordings Network link it to the Keriau Punan, with Kereho-uheng as another name. The sources also disagree on the provinces where it is spoken. Resources below are listed as found, and the identification is unresolved.

\section{Primary Educational Resources}
I found no primers, textbooks or school programmes for Kereho.

\section{Linguistic and Reference Materials}
\begin{itemize}
\item Wikipedia, ``Krio Dayak language'' (xke): an overview. \url{https://en.wikipedia.org/wiki/Krio_Dayak_language}
\item Max Planck EVA Numeral Systems database, Kereho page: number-word data. \url{https://lingweb.eva.mpg.de/channumerals/Kereho.htm}
\end{itemize}

\section{Religious and Community Resources}
\begin{itemize}
\item Joshua Project profile of the Keriau Punan: community information from a missions database. \url{https://joshuaproject.net/people_groups/14478}
\end{itemize}

\section{Digital and Audio Resources}
\begin{itemize}
\item Global Recordings Network, language page for xke: I did not check whether recordings exist. \url{https://globalrecordings.net/en/language/xke}
\end{itemize}

\section{Curriculum Development Context}
Documentation is very thin, and which community a curriculum would serve needs to be settled first. The Joshua Project source reports a shift toward Indonesian among the Keriau Punan. Any curriculum work would need community consultation and new documentation.

\end{document}